%% 抛物诱导与限制
\section[抛物诱导与限制]{抛物诱导与限制（5 月 6 日）}

\srcpage{76}

我们一直在讨论几何 Satake 定理。这个定理是关于仿射 Grassmann 簇上 $G(O)$-等变层的范畴（或 $G(O)\backslash G(K)/G(O)$ 上的层）的；它是一个带有对称幺半结构的半单阿贝尔范畴。配上到 $\Vect$ 的纤维函子，我们就会得到一个 Tannakian 范畴，从而得到某个约化群的表示范畴。这就是我们今天要讨论的内容，框架是抛物诱导与抛物限制。这是一个把 $G$ 的表示与极大环面 $T$ 的表示联系起来的框架。

设 $G$ 是复约化群，例如 $\mathrm{GL}_n(\C)$、$\mathrm{SL}_n(\C)$ 等等。我们通过最高权理论研究有限维 $G$-表示。Borel 子群 $B\subset G$ 是使得 $G/B$ 为射影簇的极大子群。选一个 Borel 子群 $B$；$G/B$ 称为 $G$ 的旗簇（flag variety）。对 $G=\mathrm{GL}_n(\C)$，$B$ 的标准取法是上三角矩阵子群。极大环面，作为 $G$ 的子群来定义，同时也是 Borel 模去一个子群 $N$ 的商（对标准 Borel，$N$ 是严格上三角矩阵子群）。于是我们有一个对应
\[
  \begin{tikzcd}
    & \mathrm{pt}/B \arrow[dl] \arrow[dr] & \\
    \mathrm{pt}/G & & \mathrm{pt}/T.
  \end{tikzcd}
\]

于是我们可以试着写下推拉映射。给定 $V\in\Rep_G$，我们可以拉回到 $\mathrm{pt}/B$，这不过是把表示限制到 $B$。对推前，我们需要多一点信息：$\mathrm{pt}/B\to\mathrm{pt}/T$ 沿 $\mathrm{pt}\to\mathrm{pt}/T$ 的拉回是 $\mathrm{pt}/N$。于是给定 $V|_B$，取它的 $N$-不变量，这现在是 $B/N=T$ 的表示。$N$-不变量就是通常所谓最高权向量，而 $T$-作用就是那个最高权。这个从 $G$-表示造出 $T$-表示的过程，并不是通常经由 $T\subset G$ 的限制方式，它称为抛物限制（parabolic restriction）。

\srcpage{77}

抛物诱导是对偶的映射。从一个不可约 $T$-表示出发，即一个特征标 $\C_\lambda$，把它拉回到 $\mathrm{pt}/B$。由于
\[
  \mathrm{pt}/B\simeq G\backslash G/B,
\]
$\C_\lambda$ 到 $\mathrm{pt}/B$ 的拉回被等同于 $G/B$ 上一个 $G$-等变向量丛。从旗簇到 $\mathrm{pt}$ 的推前，同时追踪 $G$-作用，给出一个 $G$-表示；这就是最高权表示的构造。

在物理中，这个过程定义了规范群为 $G$ 的规范场论与规范群为 $T$ 的规范场论之间的一个界面（interface）。这大致意味着映入上述对应图。例如，如果你了解幂零锥的 Springer 消解，你可以把它表述为映入那个对应图。

现在我们要把它升级为球面范畴的抛物诱导与抛物限制。上述对应图的类比是
\[
  \begin{tikzcd}
    & \Gr_B \arrow[dl] \arrow[dr] & \\
    \Gr_G & & \Gr_T.
  \end{tikzcd}
\]
而我们已经看到 $\Gr_T\cong\Lambda$，一个格。阿贝尔对偶随后等同出 $\mathrm{Sph}_T\cong\Rep_{T^\vee}$，其中 $T^\vee$ 是对偶环面。我们以前没见过 $\Gr_B$；而由于 $B$ 不是约化的（它是可解的），这里有点意思。我们还更关心 $\ul{\Gr}_B$，其中我们加入了 $B$-等变性：
\[
  \ul{\Gr}_B := B\backslash\Gr_B = B(O)\backslash B(K)/B(O).
\]
空间 $\Gr_B$ 参数化曲线 $C$ 上的主 $B$-丛，连同远离一点处的平凡化。$G$-丛的平凡化比给出到 $B$-丛的约化（那是一整组旗的数据）强得多。因此我们有一个从 $C\setminus\mathrm{pt}$ 到「参数化平凡丛中的旗」的那个东西的映射，而那是 TODO。所以仿射 Grassmann 簇上的一个点，即 $C\setminus\mathrm{pt}$ 上有平凡化的 $G$-丛，给出 $C\setminus\mathrm{pt}$ 上配的 $G/B$-丛的一个截面。而由于 $G/B$ 射影，紧合性的赋值判别法允许我们延拓到整个 $C$。这在族中不成立：它只对域值点有效。

因此映射 $\Gr_B\to\Gr_G$ 在 $k$-点上是双射。但它不是同胚——实际发生的事相当漂亮：你可以把 $\Gr_B$ 拆成空间 $\Gr_{B,\lambda}$ 的不交并，指标 $\lambda\in\Lambda=\Gr_T$，经由映射 $\Gr_B\to\Gr_T$。TODO：我想我可能误解了这里的什么。

你可以考虑 $\Gr_G$ 被 $LN=N(K)$ 作用的轨道：$N$ 正是从 $B$ 到 $T$ 的过程中被杀掉的东西。这些轨道既无穷维又有无穷余维，称为半无限轨道（semi-infinite orbit）。

这与 Floer 理论中也出现的那类现象有关：那里有无穷维 Morse 理论，把一个空间分层成无穷余维且无穷维的块。

要点是：我们有一个函子 $\mathrm{Sph}_G\to\mathrm{Sph}_T\simeq\Rep_{T^\vee}$，描述方式是取一个层 $\cF$，对 $N(K)$-轨道做 $!$-限制（shriek restriction），再取上同调。\yizhu{$!$-限制指沿开浸入取的「下叹号」拉回，与 $*$-限制相对；在字符串层论中常记为 $j^!$。}{I（术语背景）}这是跨上述对应图的推拉。于是我们得到一个 $\Lambda$-分次向量空间。忘掉分次，我们得到一个函子 $\mathrm{Sph}_G\to\Vect$，这就是纤维函子。一直飘着的额外结构，例如 $\Lambda$-分次，让人能证明纤维函子是对称幺半的；它显然也是忠实的，看轨道便知。于是我们可以跑 Tannakian 表示。

\begin{theorem}[几何 Satake]
作为对称幺半范畴，$(\mathrm{Sph}_G,*)$ 等价于 $(\Rep_{G^\vee},\otimes)$；其中 $G^\vee$ 是纤维函子的自同构群，且恰是 $G$ 的 Langlands 对偶群。
\end{theorem}

历史是曲折的：Lusztig 首先证明了类似的结果，Ginzburg 在此基础上进一步发展。Drinfeld 造出了因子化图像，这很关键；Mirković–Vilonen 把各块拼成了这个定理。

Langlands 对偶群 $G^\vee$ 是一个约化群，其极大环面是 $T^\vee$，即 $T$ 的对偶。这是因为我们造了一个函子 $\mathrm{Sph}_G\to\Rep_{T^\vee}$：Tannakian 重构告诉我们，这给出一个典范的子群 $T^\vee\subset G^\vee$：这次我们不必选择它。可以验证 $G^\vee$ 的 Weyl 群同构于 $G$ 的 Weyl 群，而 $G$ 与 $G^\vee$ 的根数据互为对偶。同样正确的是，$Z(G^\vee)$ 与 $\pi_1(G)=\pi_0(\Gr_G)$ 的 Cartier 对偶等同：在连通分支上，这个构造还原出 Cartier 对偶。

\srcpage{78}

还有一个问题：我们在什么域上工作。答案是：我们把 $G^\vee$ 作为球面范畴中层的系数域上的群得到——我们一直在 $\C$ 上工作，但并非必须如此。Mirković–Vilonen 证明了你可以取任意系数，甚至 $\Z$；而 $G^\vee$ 会定义在那个系数环 $R$ 上，因为你在 $R$ 上做 Tannakian 重构。这是件大事：非代数闭域上的约化群表示可能非常复杂。我们先前把几何 Satake 定理想成用表示论（在特征零的代数闭域上更容易）来理解球面范畴（更难）；但在别的系数环上（例如模特征情形），定理仍然成立，即便范畴不再是半单的。人们用球面范畴来推进对 $\Rep_G$ 的理解。流向反过来了！

\begin{remark}
「任意系数？」那 $\mathbb{S}$ 上呢？球面范畴在同伦论中仍有意义，那里我们只取谱的层；但有一个大问题：这个范畴不再是对称幺半的：它只是 $E_3$，而在导出设定中 $E_3$ 不是 $E_\infty$。导出几何 Satake 定理能做什么是已知的，但许多东西很怪：在 $\mathbb{S}$ 上你甚至无法定义好的行列式映射，从而无法定义 $\mathrm{SL}_n$。\yizhu{$\mathbb{S}$ 指球谱（sphere spectrum），即整数环在谱意义下的类比；「$\mathbb{S}$ 上的 $\mathrm{SL}_n$」指把群概形改在球谱上定义。}{I（术语背景）}
\end{remark}
